\begin{afterword}
\summaryhead

The post names a class of ``wrong questions'': questions that cannot be answered on their own terms and can only be dissolved, by understanding the mental process that makes them feel like questions. A good cue is that you cannot imagine any concrete state of the world that would answer the question. The author gives three examples of increasing difficulty: whether a tree falling unheard makes a sound, free will, and why anything exists at all. The author does not know the answer to the last, but predicts that it will come as an insight showing the question was wrong.

Mystery and confusion, the post says, exist in the mind, not in reality. Trying to answer such questions produces mysterious answers that predict nothing. In the author's experience, unanswerable questions are always solvable, unlike questions such as what Queen Elizabeth I thought one morning, whose answers are easy to imagine but lost.

\responsehead

The post's test is useful and its distinction is clear. A question whose answer we can imagine but cannot recover, like Queen Elizabeth's thoughts, differs from one where we cannot say what an answer would look like. The test is also old. Wittgenstein's \textsc{Tractatus} (1921) has: ``If a question can be put at all, then it can also be answered.''

The free-will example does not fit the test as well as the post assumes. The post asks what state of affairs, in deterministic or random physics, ``could ever correspond to having free will.'' Hume named one in 1748: a power of acting according to the determinations of the will, which belongs to anyone ``not a prisoner and in chains.'' That state of affairs exists in deterministic physics. The question may be unanswerable for some notion of free will that the post leaves unstated, but not for Hume's. Three months later the author gave a compatibilist answer, like Hume's.

The central claim rests on experience the reader does not see. Unanswerable questions are ``always solvable, at least in my experience,'' and the reader is assured ``that it is how these things work.'' Of the three examples, one had been treated in earlier posts, one was set that day as homework, and for the third the author says ``I don't know the answer.'' The claim that trying to answer such questions ``inevitably'' produces mysterious answers is likewise unsupported, apart from one hypothetical answer that the author offers and rejects.

In short: a useful test with an old pedigree, a free-will example that an old answer meets, and a claim that such questions are always solvable resting on experience the post does not show.
\end{afterword}
